Large language models (LLMs) have rapidly become a default interface for asking questions about the scientific
literature. A researcher, student, or clinician can now type a natural-language question and receive a fluent,
confident answer in seconds. This fluency is also the source of a serious problem. The same models that summarize
a field convincingly will, with equal confidence, assert findings that no study supports, attribute results to
papers that do not contain them, and invent citations that look entirely plausible. In settings where the cost of
a wrong answer is high --- medicine, biology, public policy, or any domain where decisions follow from evidence
--- a confident but ungrounded answer can be worse than no answer at all.

The difficulty is not only that models occasionally make mistakes. It is that a single-model answer offers the
reader no way to tell a grounded conclusion from a fabricated one. The output arrives as a finished paragraph
with no visible separation between what was retrieved, what was inferred, and what was simply generated. There is
no explicit confidence, no list of the evidence that was actually consulted, and no record of the reasoning steps
that produced the conclusion. For a tool that is meant to support research, this opacity is the central failure
mode. The question this thesis addresses is therefore not ``can an LLM answer a scientific question?'' but rather
``how can a system answer a scientific question in a way that is grounded, calibrated, and inspectable?''

This thesis presents \textbf{AI-Scientist}, a multi-agent orchestrated framework for scientific claim
verification and adaptive reasoning. Instead of relying on one model to read, reason, and conclude in a single
opaque step, AI-Scientist decomposes the task into specialized agents coordinated by an Orchestrator: a Retrieval
agent gathers relevant papers, a Claim Extraction agent turns them into discrete checkable statements, a
Verification agent weighs each statement against evidence, a Critic agent challenges weak or single-source
conclusions, a Synthesis agent composes a grounded summary, and a Report agent assembles the final
evidence-backed output with confidence scores and a reasoning trace. The design treats verification, rather than
generation, as the primary task.

\section{Research Context and Gap}

The research landscape already contains strong prior work on retrieval-augmented generation (RAG), on automated
fact-checking, and on multi-agent LLM systems. Retrieval-augmented approaches reduce hallucination by conditioning
generation on retrieved passages, and a large body of work studies how to fetch, rank, and integrate that
context. Multi-agent and ``orchestrator'' designs have shown that decomposing a task across cooperating models can
improve robustness on complex reasoning problems. AI-Scientist does not claim to be the first system to retrieve
evidence before answering, nor the first to use multiple cooperating agents. Such claims would be inaccurate and
would obscure the actual research opportunity.

The genuine gap addressed by this thesis is different. Most retrieval-augmented assistants still optimize for a
single fluent answer and treat retrieval as a silent preprocessing step; they rarely expose, per claim, whether
the evidence actually supports the statement, how strongly, from how many independent sources, and what should be
done when the evidence is thin or contradictory. In other words, the literature has established that grounding
helps, but it has not adequately framed scientific question answering as an explicit, auditable
\emph{claim-by-claim verification process} with calibrated confidence and a first-class option to abstain. This is
where AI-Scientist contributes: it adds an orchestrated verification and critique layer on top of an
already-established grounding setting, and it makes that layer the visible product rather than an implementation
detail.

\section{Thesis Position}

AI-Scientist demonstrably answers research questions end to end: given a question from any academic domain, it
retrieves evidence, extracts claims, assigns a grounded verdict and confidence to each, and produces a report in
which every conclusion is traceable to its sources. The position taken in this thesis builds on that capability
rather than overstating it. AI-Scientist is claimed as a reproducible, evidence-grounded, verification-first
research assistant that is dependable when the evidence is clear and appropriately cautious when it is not. The
main contribution is consequently not raw answer fluency or breadth of coverage; it is the engineering of
\emph{trustworthy} scientific answering --- knowing \emph{when to say ``insufficient evidence''} --- together with
an orchestration loop in which a critic can force targeted re-verification before a conclusion is treated as
settled.

The framing of the project rests on four commitments that recur throughout the thesis:

\begin{enumerate}
\item \textbf{Grounding over generation.} Every conclusion must be tied to specific retrieved evidence. A claim
with no supporting sentence in the evidence pool is reported as unsupported, not paraphrased into apparent fact.
\item \textbf{Calibrated confidence over false certainty.} The system attaches an explicit confidence to each
verdict, weighted by evidence strength, corroboration, and source independence, and it deliberately reserves the
highest confidence for genuinely well-supported claims.
\item \textbf{Critique over single-pass acceptance.} A dedicated Critic agent flags insufficient, low-confidence,
single-source, and contradicted claims, and can trigger a second retrieval-and-verification pass aimed precisely
at the weak claim.
\item \textbf{Traceability over opacity.} The pipeline records an orchestration trace and per-claim evidence so
that a reader can reconstruct exactly why the system reached each conclusion.
\end{enumerate}

These commitments are valuable precisely because they constrain the system. A central strength of AI-Scientist is
that it is designed to under-claim rather than over-claim: when evidence is missing or conflicting, it says so.

\section{Problem Motivation}

The motivation for AI-Scientist comes from the mismatch between how scientific questions are actually settled and
how single-model assistants answer them. In real research practice, a claim is not accepted because it sounds
authoritative; it is accepted because it is supported by evidence, ideally corroborated across independent
sources, and it remains open to challenge when contradictory findings appear. A useful research assistant
therefore needs at least three properties.

First, it must be \textbf{evidence-grounded}. Scientific answers cannot be trusted on the strength of phrasing
alone. Retrieved papers, their abstracts, and the specific sentences that bear on the question provide the
material against which a claim must be checked.

Second, it must be \textbf{uncertainty-aware}. A generated statement should not be treated as established merely
because it is plausible. Practical verification demands an explicit notion of confidence, an account of how many
independent sources agree, and a willingness to report uncertainty honestly.

Third, it must be \textbf{self-critical}. The system must be able to recognize when a conclusion rests on a single
source, when confidence is low, or when the evidence is contradictory, and it must be able to act on that
recognition by seeking more evidence rather than papering over the gap.

AI-Scientist is designed around these three requirements. The resulting system is not a general-purpose oracle;
it is a scoped verification assistant whose primary value lies in disciplined reasoning under noisy and incomplete
evidence.

\section{Research Questions}

This thesis is organized around the following research questions:

\begin{enumerate}
\item Can a multi-agent orchestrated framework improve scientific claim verification and reasoning traceability
compared with a direct single-agent reader and a standard RAG pipeline?
\item Does grounding answers in a layered evidence base --- user-uploaded papers, a curated corpus, and live
scholarly sources --- reduce unsupported claims relative to ungrounded answering?
\item Can per-claim confidence, weighted by evidence strength and source independence, give the user a usable
signal for when to trust a conclusion?
\item Does an explicit critic-driven re-verification loop improve the handling of weak, single-source, or
contradicted claims?
\item Can the system remain domain-general, accepting research questions across fields while still rejecting
clearly non-research queries?
\end{enumerate}

These questions keep the work grounded. They separate retrieval from verification, verification from critique, and
confidence from certainty. They also ensure that the thesis does not collapse into a vague ``AI for science''
claim.

\section{Contributions}

The contributions of this thesis are summarized as follows:

\begin{enumerate}
\item \textbf{A system contribution:} AI-Scientist, an end-to-end multi-agent framework that integrates
retrieval, claim extraction, verification, critique, synthesis, and reporting behind a single Orchestrator, with a
Streamlit interface and a programmatic API.
\item \textbf{A methodological contribution:} a verification-first formulation of scientific question answering,
in which the unit of output is a checked claim with a verdict, an explicit confidence, and attributed evidence,
rather than a single undifferentiated paragraph.
\item \textbf{An evidence-layering contribution:} a layered, source-weighted retrieval strategy that combines
user-uploaded papers, a curated offline corpus, and live results from PubMed Central and arXiv, with an explicit
priority that favours user-provided and freshly retrieved evidence.
\item \textbf{An orchestration contribution:} an iterative critic-driven re-verification loop that targets weak
claims specifically, together with a transparent orchestration trace that records how each conclusion was reached.
\item \textbf{A comparative framing:} an explicit positioning against single-agent and RAG baselines that
clarifies what the multi-agent design adds and where its limits lie.
\end{enumerate}

\section{Thesis Organization}

The remainder of this thesis is organized as follows. Chapter~2 reviews the background and related work on LLM
hallucination, retrieval-augmented generation, automated fact-checking, and multi-agent reasoning systems.
Chapter~3 defines the problem, the scope, and the requirements of AI-Scientist as a verification system.
Chapter~4 presents the system design and architecture, including the orchestrated agent pipeline and the layered
evidence model. Chapter~5 describes the implementation of the agents and their supporting modules. Chapter~6
explains the evaluation methodology, baselines, and the qualities the system is judged on. Chapter~7 reports the
evaluation through representative cross-domain case studies and a comparison with the baselines. Chapter~8
discusses the broader meaning of these findings. Chapter~9 examines threats to validity and current limitations.
Chapter~10 concludes the thesis and outlines future research directions.

\section{Chapter Conclusion}

This chapter has introduced the AI-Scientist thesis, its motivation, research questions, and contributions. It has
also established the central framing used throughout the document: answering a scientific question well is not a
generation problem but a verification problem that requires evidence grounding, calibrated confidence, explicit
critique, and traceable reasoning. The next chapter situates this framing within the related literature.
